\documentclass[12pt]{article}
\usepackage[utf8]{inputenc}
\usepackage[T1]{fontenc}
\usepackage{hyperref}
\title{Shuhi Language Resources for K-12 Curriculum Development}
\author{Compiled by Claude Sonnet 5.5 (claude-sonnet-5-5) for Andr\'as Kornai}
\date{}
\begin{document}
\maketitle

\section{Introduction}
Shuhi is the speakers' name for the language that is also called Shixing or Xumi (Chinese \emph{Shixing hua}). It is spoken by about 1,800 people along the Shuiluo River in Muli Tibetan Autonomous County, Sichuan, China. Katia Chirkova describes two varieties, Upper Xumi and Lower Xumi, with limited mutual intelligibility. Classifications differ: Wikipedia calls it Qiangic, while SIL-based sources call it Naic. Search found scholarly descriptions only, and no school materials.

\section{Primary Educational Resources}
No primers, textbooks, or school programmes were found in search results.

\section{Linguistic and Reference Materials}
\begin{itemize}
\item Wikipedia article on the Shixing language. \url{https://en.wikipedia.org/wiki/Shixing_language}
\item Chirkova, Katia, \emph{Lower Xumi} (2012, HAL open-access PDF). A description of the Lower Xumi variety. I did not read it. \url{https://halshs.archives-ouvertes.fr/hal-01136684v1/file/2012\%20Lower\%20Xumi\%20HAL.pdf}
\item Glottolog reference entry for Chirkova (2014). \url{https://glottolog.org/resource/reference/id/319917}
\item ASJP word list for Shixing. \url{https://asjp.clld.org/languages/SHIXING}
\item Cocoon (CNRS/Huma-Num) archive listing for code sxg. \url{https://cocoon.huma-num.fr/exist/crdo/lexvo.xql?code=sxg}
\end{itemize}
Search summaries also name, without a URL I saw: Chirkova's grammatical sketch ``Sh\v{i}x\=ing, a Sino-Tibetan language of South-West China'' (\emph{Linguistics of the Tibeto-Burman Area} 32, 2009), and \emph{Muli Xumi fangyan} (Minzu Chubanshe, 2013). They are listed for follow-up only.

\section{Religious and Community Resources}
\begin{itemize}
\item Asia Harvest people-group profile of the Shixing. It is a mission-oriented demographic profile. \url{https://people-groups.asiaharvest.org/people-group-profiles/sichuan/shixing.pdf}
\end{itemize}
No scripture translations or local community organizations were found.

\section{Digital and Audio Resources}
\begin{itemize}
\item Endangered Languages Project page for Shixing (Xumi), with a bibliography. \url{https://endangeredlanguages.com/lang/5431}
\end{itemize}
No audio, video, or software were found beyond the above.

\section{Curriculum Development Context}
Shuhi is still spoken by most adults, but not all young people. No orthography, dictionary, or classroom materials were found in search results. Chirkova's published work is the main documentation and could underlie a community orthography and a first wordlist. School instruction in Muli is in Chinese, so any Shuhi-medium content would have to begin as a supplementary heritage programme developed with the community. Local contacts and programmes were not found and should not be assumed.
\end{document}